\chapter{Introduction}

A \emph{perceptron}\keyword{Perceptron} is a fundamental computational unit of neural networks, inspired by the behavior of biological neurons. It receives multiple inputs, computes a weighted combination of those inputs, and produces an output after applying a nonlinear activation function.

Formally, let $\vect{x} \in \R^n$ denote an input vector. A perceptron computes
\[
    y = \sigma\left(\vect{w}^\top \vect{x} + b\right),
\]
where $\vect{w} \in \R^n$ is a vector of \emph{weights}, $b \in \R$ is a \emph{bias term}, and $\sigma(\cdot)$ is an \emph{activation function}.

\begin{figure}[tb]
    \centering
    \begin{tikzpicture}[node distance=12mm and 18mm, font=\small]

% Inputs
\node (x1) {$x_1$};
\node (x2) [below=of x1] {$x_2$};
\node (vdots) [below=of x2] {$\vdots$};
\node (xn) [below=of vdots] {$x_n$};

% Summing junction
\node (sum) [sum, right=35mm of x2] {$+$};

% Bias input
\node (b) [above=18mm of sum] {$b$};

% Activation block
\node (act) [block, right=20mm of sum] {$\sigma(\cdot)$};

% Output
\node (y) [right=18mm of act] {$y$};

% Input arrows
\draw[arrow] (x1) -- node[above, sloped] {$w_1$} (sum);
\draw[arrow] (x2) -- node[above, sloped] {$w_2$} (sum);
\draw[arrow] (vdots) -- (sum);
\draw[arrow] (xn) -- node[below, sloped] {$w_n$} (sum);

% Bias arrow
\draw[arrow] (b) -- (sum);

% Sum to activation
\draw[arrow] (sum) -- node[above] {$\vect{w}^\top\vect{x} + b$} (act);

% Activation to output
\draw[arrow] (act) -- (y);

\end{tikzpicture}
    \caption{A perceptron receives input features, computes a weighted sum with bias, and applies an activation function.}
\end{figure}

The quantity $\vect{w}^\top \vect{x} + b$ defines a linear function of the input. The set of points satisfying
\[
    \vect{w}^\top \vect{x} + b = 0
\]
forms a hyperplane in $\R^n$. This hyperplane divides the input space into two half-spaces, corresponding to the two possible outputs of the perceptron.

In a binary classification setting, the perceptron assigns a label according to the sign of this linear function:
\[
    y =
    \begin{cases}
        1, & \text{if } \vect{w}^\top \vect{x} + b > 0, \\
        0, & \text{otherwise.}
    \end{cases}
\]

Although a single perceptron can represent only linearly separable decision boundaries, networks composed of multiple perceptrons can approximate highly complex nonlinear functions.
